\paragraph{Conjugate-coordinate class times a fuzzy factor: fibre spectrum, proper-time interpolation between the two readings, and the background torsion (conditional; bounded no-go; correction of a regime argument of the preceding paragraph).}
This paragraph supplements the preceding one, in which the interpolation between the fibre trace and the averaged-frame reading, the fibre spectrum and the fibre-flux part were not computed. Keep its setting (Poisson limit on \(M\), leaf \(\mathbb{CP}^{1,2}\), \(T^a=M^{a4}\), \(y^\mu=M^{\mu5}\), \(m=M^{45}\), units \(R=1\), the split \(d=4\), the level-one amplitudes, the leading symbol at a frozen base point, the orientation-even mixed term, a formal evaluation with the Lorentzian \(\eta\), inert amplitudes set to zero), write \(t^\mu=T^\mu\), \(\coth\tau=y^0/y^4\) at \(y^i=0\), and let the loop modes be \(\phi(y)Y_{Jm}(n)\) with fibre harmonics of order \(J\le n\). The scale of \(K\) is measured by the eigenvalues \(M\) of the \(K\)-part of the loop operator, in the units of \(\Box\).

\emph{Fibre spectrum.} In the polynomial Lie--Poisson algebra, for both signs of the bracket, \(\Box t^\mu=-3t^\mu\) and \(\Box(t^\mu t^\nu)=-6t^\mu t^\nu+2\eta_{ab}M^{a\mu}M^{b\nu}\), and on the orbit \(\eta_{ab}M^{a\mu}M^{b\nu}=c_\theta\eta^{\mu\nu}+y^\mu y^\nu-t^\mu t^\nu\) (\(40\) random points, residual below \(10^{-15}\)). On the fibre over a base point \(t\cdot y=0\) and \(t\cdot t\) is constant, and for \(C\) symmetric, transverse to \(y\) and traceless on \(y^\perp\), \(\Box(C\cdot t^J)=-J(J+2)\,C\cdot t^J\) (\(J=1,2,3\), two base points, one with \(y^i\neq0\), relative residual below \(10^{-15}\)). The fibre levels of the loop operator \(-\Box\) at zero base momentum are therefore \(\mu_J=J(J+2)\); they agree with the value \(3J\) of Equation~(7.5) of~\cite{Steinacker2023} at \(J=1\) and exceed it for \(J\ge2\), and the two are not reconciled here. The cross term \(2\eta_{ab}\{T^a,\phi\}\{T^b,Y\}\) for a plane wave \(\phi\) of momentum \(q\) maps each level \(J\) into itself, with singular values \((2m/\lvert y^0\rvert)\bigl((y^4q_0J)^2+(\ell\lvert\mathbf q\rvert m_L)^2\bigr)^{1/2}\), \(m_L=-J,\dots,J\), at \(y^i=0\) (Monte Carlo agreement within four per cent for \(J=1,2\)); it is comparable with \(m^2q^2\) and \(\mu_J\) when \(m\lvert q\rvert\sim J\) and is dropped below, as in the preceding paragraph.

\emph{Proper-time weights.} A vertex that multiplies by a fibre function \(f=f_0+f_2\) (harmonic orders \(0\) and \(2\)) connects \(J\) to \(J\) and \(J\pm2\), with \(\sum_{m,m'}\langle Jm|f|J'm'\rangle\langle J'm'|g|Jm\rangle=(2J+1)\bigl[\delta_{JJ'}f_0g_0+c(J,J')\langle f_2g_2\rangle\bigr]\), \(c(J,J')=\langle J0\,20|J'0\rangle^2\). At second order the intermediate level enters through \(e^{-s\Delta a}\), \(\Delta=\mu_{J'}-\mu_J\), \(0\le a\le1\). The torsion vertex is already of order \(q^2\), so its weight is \(\int_0^1e^{-s\Delta a}\,\mathrm da\). For the scalar operator \(-\partial\gamma\partial\) with cometric wave \(h\), \(G=\bar\gamma^{-1}\), the part of order \(q^2\) is proportional to \(\int_0^1e^{-s\Delta a}[a(1-a)F_1+(1-2a)^2F_2]\,\mathrm da\) with \(F_1=-\tfrac14(q\bar\gamma q)\bigl((\operatorname{tr}hG)^2+2\operatorname{tr}(hGhG)\bigr)-(qhq)\operatorname{tr}(hG)\) and \(F_2=\tfrac12\,qhGhq\) (exact Gaussian moments, checked independently in both signatures); at \(\Delta=0\) this gives the scalar density of the preceding paragraph, which equals \(\tfrac14(\tfrac16F_1+\tfrac13F_2)/\sqrt{\lvert\det\bar\gamma\rvert}\) exactly. Hence, with level weights \(D(s)=\sum_J(2J+1)e^{-s\mu_J}\) on the averaged parts and the corresponding sums \(B_i(s)\) on the second-harmonic parts, the index form is \(W_X^{(0)}\,\mathrm{IX}_1+W_X^{(2)}(\mathrm{IX}_2-\mathrm{IX}_1)\) and the scalar form is \(W_R^{(0)}\,\mathrm{SC}_1+W_R^{(1)}L_1+W_R^{(2)}L_2\), where \(\mathrm{IX}_{1,2}\), \(\mathrm{SC}_{1,2}\) are the averaged-frame and fibre-trace forms of the preceding paragraph, \(L_1+L_2=\mathrm{SC}_2-\mathrm{SC}_1\) (exact at the seven sampled points), \(L_2\) vanishes on the five traceless polarisations, and each \(W\) is the integral of \(s^{-2}\tau^K(s)\) or \(s^{-2}(v^K-\tau^K)(s)\) against the corresponding level weight. The cone carries only the weight \(D\), so the mixed-term part of the fixed-state gap on the generalised eigendirections of the five traceless polarisations is \(-W_X^{(2)}(\lambda_i+2)\), with \(\lambda_i\) the eigenvalues of \(\mathrm{IX}_2\) relative to \(\mathrm{EH}(\bar h)\). The form therefore interpolates exactly between the two readings, and a single coefficient on the five holds only where \(W_X^{(2)}=W_R^{(1)}=0\).

\emph{Which regime.} For \(\mu_J\gg M\) a level contributes \((2J+1)\int\mathrm ds\,s^{-2}\tau^K(s)e^{-s\mu_J}\simeq(2J+1)\tau_2/\mu_J\), with \(\tau^K\simeq\tau_2s^2\) at small \(s\): the levels are suppressed by a power, not exponentially, and the transitions \(J\to J\) of the second harmonic have \(\Delta=0\). The ratio \(\beta_X=W_X^{(2)}/W_X^{(0)}\) is small only if \(M\sum_{J\ge1}(2J+1)/\mu_J\ll1\) and \(M\ll\mu_2\), and tends to one as \(n\to\infty\) at fixed \(M\), the sum growing as \(2\ln n\). For \(\tau^K\propto s^2e^{-sM}\) and \(n=30\), \(\beta_X=0.065\), \(0.40\), \(0.86\), \(0.96\) at \(M=0.01\), \(0.1\), \(1\), \(8\) (\(0.19\), \(0.69\), \(0.94\), \(0.97\) for the levels \(3J\)), and at the base point \((y^0,y^4)=(5,3)\) with \(q=(\tfrac37,\tfrac12,-\tfrac23,\tfrac15)\) the half-spread of the five eigenvalues of the index form relative to \(\mathrm{EH}(\bar h)\) is \(2.42\,\beta_X\) about the averaged-frame value \(-2\) (numerical). The argument of the preceding paragraph that the full trace approaches the averaged-frame forms when the fibre gaps exceed the scale of \(K\) is therefore incomplete: the condition is \(M\sum_{J\ge1}(2J+1)/\mu_J\ll1\), which, at fixed \(M\), fails as \(n\to\infty\), since the sum grows as \(2\ln n\).

\emph{Background torsion.} The flux operator of a fibre-constant loop mode is \(\mathrm ik_\mu T_{ab}{}^\mu\) with \(T=\bar T+\delta T_0\), \(\bar T^{ab\mu}=\eta^{a\mu}y^b-\eta^{b\mu}y^a\) and \(\delta T_0\) the fibre average of the level-one torsion, and at frozen coefficients the index part of a species is the Gaussian average over \(k\) of \(\operatorname{tr}e^{-sZ(k)}-\dim\), \(Z=-\mathrm i\Sigma^{ab}T_{ab}{}^\mu k_\mu\), exactly. Its expansion \(sQ_1+s^2Q_2+\dots\) gives for the Dirac representation \(Q_1=m^4\,\mathrm{IX}_1\) to machine precision at ten sampled points, so on these modes the index density is the torsion invariant of the averaged frame, and the vector representation gives twice it. The part of \(Q_2\) quadratic in \(\delta T_0\) is the first background-torsion correction of the mixed term (fluxes of order four on \(M\) and two on \(K\)); the same part of \(Q_2^V-2Q_2^D\) is the first term \(s_0^K(v^M-2\sigma^M)\) at order four, with the weight \(\int\mathrm ds\,s^{-1}s_0^K(s)\), logarithmic at both ends and not suppressed by the flux of \(K\). On the five traceless polarisations the forms \(\mathrm{IX}_2-\mathrm{IX}_1\), \(Q_2^D/m^4\) and \((Q_2^V-2Q_2^D)/m^4\) lie in the span of \(\mathrm{EH}(\bar h)\) and one further form, with \(Q_2^V-2Q_2^D=-6Q_2^D\) modulo \(\mathrm{EH}(\bar h)\) and \(\mathrm{IX}_2-\mathrm{IX}_1=\lambda Q_2^D/m^4\) modulo \(\mathrm{EH}(\bar h)\), \(-3.31\le\lambda\le-0.32\) (residual below \(10^{-11}\) at seven sampled points); the eigenvalues of \(Q_2^D/m^4\) relative to \(\mathrm{EH}(\bar h)\) there are one simple and two double values, one of which is \(\tfrac32\coth^2\tau\) at every sampled point. The scalar second-harmonic form \(L_1\) is not in that span. The relative size of the correction is \(s\coth^2\tau\), with \(\coth^2\tau\ge1\) on the leaf, weighted by the scale of \(K\) through \(\langle s\rangle\simeq1/M\).

\emph{Consequence.} Where the second-harmonic weight is small, \(M\sum_{J\ge1}(2J+1)/\mu_J\ll1\), the torsion parameter is \(\coth^2\tau/M\gg\coth^2\tau\sum_{J\ge1}(2J+1)/\mu_J\ge1\), so the forms of order two in the flux of \(M\), on which the averaged-frame coefficient \(c_T=\tfrac16W_R-2W_X\) and the equality \(c_A=c_T\) on the five traceless polarisations rest, are not the leading result there; where the torsion expansion holds, \(M\gg\coth^2\tau\), the weight \(\beta_X\) is of order one (it tends to one as \(n\) grows) and the fibre trace of the preceding paragraph is the leading reading. At first order in the torsion, and with the first-term part weighted by the ratio \(L/c\) of its proper-time logarithm to the flux normalisation \(\tau_2=cM^2\) of \(K\) (a parameter, not computed), the coefficient of the non-Einstein form of the index sector changes sign at one value \(M^*\); at the base point \((5,3)\) and the wave vector above, for \(2\le n\le300\), it lies at \(0.21\le M^*\le0.67\) with \(\coth^2\tau/M^*\ge4.2\) for \(L/c=0\), and at \(1.1\le M^*\le30\) with \(\beta_X\) between \(0.52\) and \(0.97\) for \(L/c=1\) and \(10\), where the cancellation is between the fibre-trace part and the first-term torsion part. The logarithm \(L\) comes from fibre-constant modes without \(K\) flux, for which the torsion expansion is not controlled at large proper time. A single coefficient there would further require the scalar coefficient of the independent form \(L_1\) to vanish and the cone at the same order to match, which is not computed. The part of the first term with two fibre fluxes, the mixed \(S^2\)--\(M\) term of Section~7.1.2 of~\cite{Steinacker2023}, has per level a weight of order \(J^2/\mu_J\) by power counting and is not computed. Within these computations no open range of the scale of \(K\) gives a single Einstein coefficient on the five traceless polarisations with \(c_A=c_T\); the range \(M\sim\coth^2\tau\), where the first-order torsion parameter is of order one, is not computed, and isolated values of the data of \(K\) at which two independent forms vanish and the cone matches are not excluded, and would lie where the fibre trace, not the averaged frame, dominates.

These statements are for the hypotheses of the preceding paragraph, fibre-constant loop modes for the torsion correction, the base-fibre cross term dropped, commutative Clebsch--Gordan coefficients with truncation at \(n\), and the single-scale model of \(\tau^K\) for the numbers quoted; the spectrum identity is exact, the forms are exact or numerical as stated. They do not compute the fibre-flux part, the higher orders in the background torsion, the cone corrections at those orders, the elimination of the inert amplitudes or the trace polarisation, do not construct \(K\), and do not identify \(\bar h\) with the emergent metric of Chapter~\ref{ch:vierbein-emergence} or bear on Assumption~\ref{ass:lorentzian}; the replicas are those of Section~\ref{sec:status-B}. Problems~B and~F remain open and PROGRAMME-CONDITIONAL.
